\clearpage
\SourceLeaf{587}\pdfbookmark[0]{Besprechungen fremder Schriften.}{anzeigen-nicht-eigner}
\thispagestyle{empty}
\vspace*{35mm}
\begin{center}
{\LARGE BESPRECHUNGEN}\par\vspace{16mm}
{\large FREMDER}\par\vspace{16mm}
{\LARGE SCHRIFTEN.}
\end{center}
\clearpage
\SourceLeaf{589}\fancyhead[L]{\small BESPRECHUNGEN FREMDER SCHRIFTEN.}
\pdfbookmark[1]{Über die Daltonsche Theorie. 1830.}{daltonsche-theorie}
\begin{center}\small Göttingische gelehrte Anzeigen.\quad 1830. Dezember 11.\end{center}

\textit{Über die \textup{\textsc{Dalton}}sche Theorie, von \textup{J. F.\ \textsc{Benzenberg}.} Düsseldorf, \textup{1830.} Bei J. E.\ Schaub.} (192 Seiten in 8°, dazu drei Steindrucktafeln).

Die von \textsc{Dalton} aufgestellte Hypothese, dass die verschiedenen Gase, aus denen die atmosphärische Luft besteht, überhaupt keinen Druck aufeinander ausüben, sondern ebenso viele gleichsam voneinander unabhängige Atmosphären bilden, hat bisher bei wenigen Physikern Zustimmung gefunden. Unter ihnen zeichnet sich aber Herr \textsc{Benzenberg} durch den unermüdlichen, warmen Eifer, mit dem er jene Hypothese seit beinahe zwanzig Jahren verteidigt, ganz besonders aus. Insbesondere hat er in der trigonometrisch-barometrischen Messung der Höhe des Monte Gregorio durch \textsc{Daubuisson} einen wichtigen Grund \textit{für} die \textsc{Dalton}sche Hypothese gefunden. Es ist klar, dass die barometrischen Höhenmessungen, wenn die \textsc{Dalton}sche Hypothese zutrifft, anders berechnet werden müssen als nach der gewöhnlichen Theorie. Beim 5260 Fuß hohen Monte Gregorio fand Herr \textsc{Benzenberg} das Ergebnis der ersten Rechnung um 16 Fuß kleiner als das der anderen. Um nahezu denselben Betrag übertraf das Ergebnis der zweiten Rechnung das der trigonometrischen Messung. Diese Differenz würde also nach Herrn \textsc{Benzenberg}s Rechnung durch Annahme der \textsc{Dalton}schen Hypothese fast vollständig beseitigt. Herr \textsc{Benzenberg} hat diese Rechnungen zuerst 1812 in \textsc{Gilbert}s Annalen der Physik veröffentlicht und sie danach
\SourceLeaf{590}wiederholt an anderen Stellen vorgetragen. Auch zu anderen Bereichen der Physik, die mit der \textsc{Dalton}schen Auffassung in Berührung kommen, etwa der Akustik und Eudiometrie, hat sich Herr \textsc{Benzenberg} geäußert. Dabei ging es ihm weniger darum, dort Gründe \textit{für} jene Hypothese zu suchen, als vielmehr darum, die Gründe zu bekämpfen, die man daraus gegen sie ableiten kann und abgeleitet hat. In der vorliegenden Schrift ist all dies erneut zusammenhängend dargestellt.

Die Anerkennung. die diese Bemühungen bei den Physikern und Mathematikern gefunden haben, scheint zumindest bei den ausländischen nicht ausreichend gewesen zu sein. Herr \textsc{Benzenberg} erzählt im Vorwort seines Buchs, dass er bei seinem Aufenthalt in Paris im Jahr 1815 \textsc{Gilbert}s Annalen in der Bibliothek des Instituts gar nicht und das Exemplar in der großen Königlichen Bibliothek unaufgeschnitten gefunden habe. Mit \textsc{Laplace} habe er über seine Theorie der Berechnung der Barometerhöhen unter \textsc{Dalton}s Hypothese gesprochen, doch dieser große Geometer sei damals schon alt gewesen.

Den deutschen Physikern kann man jedoch nicht mit Recht vorwerfen, dass sie die Sache nicht beachtet hätten. In der neuen Ausgabe des physikalischen Wörterbuchs lesen wir im Artikel Atmosphäre, dass „Herr \textsc{Benzenberg} der bedeutendste, gründlichste und eifrigste Verteidiger der \textsc{Dalton}schen Theorie ist, dass sich mit ihr die Höhen genauer berechnen lassen als ohne sie, und dass darin ein gewichtiger Beweis für ihre Richtigkeit liege.“

Aber der Verfasser dieses Artikels scheint die Rechnungen \textsc{Benzenberg}s nicht \textit{geprüft} zu haben. Alle Zahlen sind einfach aus \textsc{Gilbert}s Annalen übernommen. Dasselbe gilt für das, was über diese Rechnungen im Artikel Höhenmessung des erwähnten Wörterbuchs gesagt wird. Vielleicht hielten die Verfasser beider Artikel eine Prüfung deshalb für weniger wichtig, weil sie den barometrischen Höhenmessungen im Allgemeinen weniger Zuverlässigkeit beimaßen als Herr \textsc{Benzenberg} und die Beweiskraft seiner Rechnungen daher ohnehin nicht anerkannten. Freilich nimmt der Verfasser des ersten Artikels damit das oben Angeführte zum Teil wieder zurück.

Es kann nicht die Absicht dieser Besprechung sein, unsere eigene Ansicht darüber darzulegen, ob die \textsc{Dalton}sche Hypothese selbst zulässig ist. Wir beschränken uns vielmehr auf das, was Herr
\SourceLeaf{591}\textsc{Benzenberg} zu ihrer Unterstützung in der vorliegenden Schrift erneut vorgetragen hat, und insbesondere auf seine Berechnung der barometrischen Höhenmessungen.

Vor allem schien es nötig, zuerst zu prüfen, ob diese Rechnung \textsc{Benzenberg}s selbst richtig ist. Zu unserer Verwunderung zeigte sich dabei, dass die Rechnung unrichtig ist und eine richtig geführte Rechnung ein ganz entgegengesetztes Ergebnis liefert.

Wir müssen dies hier etwas ausführlicher und so weit nachweisen, dass jeder selbst urteilen kann, zumal es sich hier weniger um einen Rechenfehler im eigentlichen Sinn als um einen Fehler in der Argumentation handelt.

Auf der zweiten und dritten Seite des Buchs finden wir zwei tabellarische Übersichten mit je drei Spalten, die wir einzeln durchgehen müssen. (Dieselben Tabellen stehen in \textsc{Gilbert}s Annalen der Physik, Bd.\ 42, S.\ 163 und 164.)

Die erste Spalte gibt nach \textsc{Dalton} die in 100 Teilen trockener atmosphärischer Luft enthaltenen Mengen an:
\begin{center}
\begin{tabular}{lclrl}
Stickstoff & .\quad .\quad . & zu & $78.93$ & Teilen\\
Sauerstoff & . & — & $21.00$ & —\\
Kohlendioxid & & — & $0.07$ & —
\end{tabular}
\end{center}
und zwar sämtlich nach dem Volumen. (Anstelle der zweiten Zahl steht im Buch selbst $21.90$, was ein offensichtlicher Druckfehler ist.)

Die zweite Spalte gibt die spezifischen Gewichte an
\begin{center}
\begin{tabular}{lcr}
des Stickstoffs & .\quad .\quad . & $0.9691$\\
des Sauerstoffs & .\quad . & $1.1148$\\
des Kohlendioxids & & $1.5000$
\end{tabular}
\end{center}
wobei das der gewöhnlichen atmosphärischen Luft als Einheit angenommen wird. Die erste und die dritte Zahl sind \textsc{Biot} entnommen. Die zweite hat Herr \textsc{Benzenberg}, wie er selbst sagt, so berechnet, „dass das Gemisch genau das von \textsc{Biot} angegebene Gewicht trockener atmosphärischer Luft habe, nämlich $\frac{1}{10493}$ des Gewichts des Quecksilbers bei $0^0$ Réaumur und 28 Zoll Druck.“ Man sieht, dass Herr \textsc{Benzenberg} sich selbst nicht klargemacht hat, was er hier eigentlich berechnen wollte, denn das soeben Angeführte gehört nicht zu \textit{dieser} Rechnung. Offenbar kam es nur darauf an,
\SourceLeaf{592}dem Gemisch das spezifische Gewicht 1 zu geben. Im Übrigen hat diese Unklarheit hier keinen weiteren Einfluss. Die Rechnung ist aber nicht sehr genau geführt, denn aus den angegebenen Daten folgt für das spezifische Gewicht des Sauerstoffgases nicht $1.1148$, sondern $1.11447$. Dieser Fehler. ist jedoch ganz unerheblich.

Die dritte Spalte gibt die \textit{dem Gewicht nach} in 100 Teilen trockener Luft enthaltenen Teile der einzelnen Gase an, nämlich
\begin{center}\begin{tabular}{lcrl}
Stickstoff & .\quad .\quad . & $76.49$ & Teile\\
Sauerstoff & . & $23.41$ & —\\
Kohlendioxid & & $0.10$ & —
\end{tabular}\end{center}
Diese Zahlen sind offenbar nur die Produkte der Zahlen der ersten Spalte mit den jeweils zugehörigen Zahlen der zweiten.

Die vierte Spalte gibt die Höhe an, auf der jede einzelne Atmosphäre das Barometer halten würde (unter der Voraussetzung, dass die \textsc{Dalton}sche Hypothese richtig ist), oder den Anteil am Gesamtdruck, wobei dieser für trockene Luft mit $27.76$ Zoll Quecksilberhöhe angesetzt wird. Die Zahlen von Herrn \textsc{Benzenberg} lauten
\begin{center}\begin{tabular}{llcrl}
für die & Stickstoffatmosphäre & .\quad .\quad . & $21.2326$ & Zoll\\
—\quad — & Sauerstoffatmosphäre & . & $6.4986$ & —\\
—\quad — & Kohlendioxidatmosphäre & & $0.0278$ & —
\end{tabular}\end{center}
Wir kommen gleich auf die Berechnung dieser Spalte zurück.

Die fünfte Spalte enthält die spezifischen Gewichte der Gase im Vergleich zu Quecksilber beim Gefrierpunkt und bei 28 Zoll Quecksilberdruck. Diese Zahlen, nämlich der Reihe nach $\frac{1}{10830}$, $\frac{1}{9414}$, $\frac{1}{6997}$, sind nichts anderes als die Produkte der Zahlen der zweiten Spalte mit $\frac{1}{10495}$.

\Needspace{20\baselineskip}
Die sechste Spalte trägt nur die Überschrift \textit{Beständige Zahl.} Man sieht aber, dass die sogenannte Subtangente gemeint ist, also die Höhe, die eine gedachte Atmosphäre mit gleichförmiger Dichte, und zwar mit derselben Dichte wie die wirkliche Atmosphäre am Boden, haben müsste, um denselben Druck auszuüben wie diese. Für gewöhnliche trockene Luft ist diese Zahl also das Produkt aus 10495 und 28 Zoll oder $24488\frac{1}{3}$ Fuß. Für die drei einzelnen Atmosphären nach \textsc{Dalton}s Vorstellung sind die Zahlen ebenso die Produkte aus 28 Zoll und den Nennern der Brüche in der fünften Spalte. Einfacher erhält man sie, indem man $24488.33$ Fuß durch die Zahlen der zweiten Spalte dividiert. Wir führen diese Zahlen hier an,
\SourceLeaf{593}sowohl so, wie Herr \textsc{Benzenberg} sie angibt, als auch so, wie sie aus einer genaueren Rechnung folgen
\begin{center}\begin{tabular}{lcrlrl}
& & \multicolumn{2}{c}{nach Hr.\ \textsc{Benz}.} & \multicolumn{2}{c}{nach genauerer R.}\\
Stickstoff & .\quad . & $25270$ & Fuß & $25269.15$ & Fuß\\
Sauerstoff & . & $21966$ & — & $21973.01$ & —\\
Kohlendioxid & & $16326$ & — & $16325.56$ & —
\end{tabular}\end{center}

Alle bisherigen Einwände gegen Herrn \textsc{Benzenberg} sind völlig unerheblich. Wir kommen jetzt zum wesentlichen Punkt, der Berechnung der vierten Spalte. Herr \textsc{Benzenberg} erklärt überhaupt nicht, \textit{wie} er diese Rechnung ausgeführt hat. Er sagt nur, es seien Beispiele aus der Gesellschaftsrechnung. Man erkennt aber leicht, dass er die Zahlen der vierten Spalte einfach proportional zu denen der dritten angesetzt oder sie durch Multiplikation von $27.76$ Zoll mit
\begin{center}\begin{tabular}{rl}
$0.7649$ & für Stickstoff\\
$0.2341$ & für Sauerstoff\\
$0.0010$ & für Kohlendioxid
\end{tabular}\end{center}
gewonnen hat.

Und das ist falsch.

Denn nach \textsc{Dalton}s Hypothese steht der gesamte Druck der Stickstoffatmosphäre zum gesamten Druck der Sauerstoffatmosphäre im \textit{zusammengesetzten} Verhältnis der Gewichtsanteile, die diese Gase am untersten Kubikfuß gemischter Luft haben, und der Subtangenten, die den beiden Gasen zukommen. Die Gewichtsanteile allein ergeben das Verhältnis also nicht. Die Drücke müssen demnach den \textit{Produkten} aus den Zahlen der dritten und sechsten Spalte proportional sein oder, was dasselbe ist, den Quotienten, die man erhält, wenn man die Zahlen der dritten Spalte durch die der zweiten teilt, also einfach den Zahlen der ersten Spalte.

Bei einigem Nachdenken wird das auch von selbst klar, denn die Bedeutung der Zahlen der ersten Spalte lässt sich auch so ausdrücken. Der Stickstoff, der in einem Volumen von 100 Teilen gewöhnlicher trockener Luft am Boden der Atmosphäre enthalten ist, würde für sich allein bei demselben Druck wie diese Luft nur den Raum von $76.93$ Teilen einnehmen, und ebenso
\SourceLeaf{594}verhält es sich mit den anderen Gasen. Da sich also jedes dieser drei Gase nun über den Raum von 100 Teilen verteilt und als unabhängig von den übrigen gedacht wird, stehen die Quecksilberdrücke, mit denen die Gase jeweils im Gleichgewicht stehen, in folgendem Verhältnis
\begin{center}\begin{tabular}{rl}
$78.93$ & für Stickstoff\\
$21.00$ & für Sauerstoff\\
$0.07$ & für Kohlendioxid.
\end{tabular}\end{center}

\Needspace{9\baselineskip}
Man sieht also, dass Herr \textsc{Benzenberg} in seiner vierten Spalte den Gesamtdruck von $27.76$ Zoll falsch aufgeteilt hat. Die richtigen Zahlen lauten
\begin{center}\begin{tabular}{rlll}
$21.9110$ & Zoll & für & Stickstoff\\
$5.8296$ & — & — & Sauerstoff\\
$0.0194$ & — & — & Kohlendioxid.
\end{tabular}\end{center}
Die Zahlen der vierten und sechsten Spalte sind aber die Größen, aus denen der Druck der Atmosphäre in jeder Höhe nach der \textsc{Dalton}schen Hypothese berechnet werden muss. Mit den berichtigten Werten ergeben sich daher natürlich andere Ergebnisse.

Wir haben uns die Mühe gemacht, diese Rechnung für einige Höhen auszuführen. Die folgende Tabelle enthält die Ergebnisse.
\begin{center}\begin{tabular}{r|r|r|r}
& \multicolumn{3}{c}{Barometerhöhe in Zoll}\\
\multicolumn{1}{c|}{Höhe} & \multicolumn{1}{c|}{nach gewöhnl.} & \multicolumn{2}{c}{nach \textsc{Dalton}s Hypothese}\\
\multicolumn{1}{c|}{Fuß} & \multicolumn{1}{c|}{Theorie} & \multicolumn{1}{c|}{unsere R.} & \multicolumn{1}{c}{H.\ \textsc{Benz}.\ R.}\\\hline
$5000$ & $22.6332$ & $22.6350$ & $22.6179$\\
$10000$ & $18.4532$ & $18.4589$ & $18.4314$\\
$15000$ & $15.0452$ & $15.0555$ & $15.0221$\\
$20000$ & $12.2666$ & $12.2814$ & $12.2458$
\end{tabular}\end{center}

Für die Barometerhöhe nach der gewöhnlichen Theorie haben wir hier die Zahlen unserer eigenen Rechnung angesetzt. Die Werte, die Herr \textsc{Benzenberg} für dieselben Höhen ebenfalls nach der gewöhnlichen Theorie berechnet hat (S.\ 12), weichen zum Teil etwas davon ab. Für die Richtigkeit unserer Rechnung können wir einstehen.

Die Barometerhöhe nach der \textsc{Dalton}schen Hypothese weicht also um folgende Beträge von der Barometerhöhe nach der gewöhnlichen Theorie ab.
\SourceLeaf{595}\begin{center}\begin{tabular}{rl|rl|rl}
& & \multicolumn{4}{c}{Unterschied der Barometerhöhe}\\
\multicolumn{2}{c|}{Höhe} & \multicolumn{2}{c|}{nach unserer R.} & \multicolumn{2}{c}{nach Hr.\ \textsc{Benz}.}\\\hline
$5000$ & Fuß & $+0.0018$ & Zoll & $-0.0153$ & Zoll\\
$10000$ & — & $+0.0057$ & — & $-0.0218$ & —\\
$15000$ & — & $+0.0103$ & — & $-0.0231$ & —\\
$20000$ & — & $+0.0148$ & — & $-0.0208$ & —
\end{tabular}\end{center}

Einer bestimmten Höhe entspricht also nach \textsc{Dalton}s Hypothese ein höherer Barometerstand als nach der gewöhnlichen Theorie. Herr \textsc{Benzenberg} nimmt dagegen einen niedrigeren an. Entsprechend ergibt sich aus einem bestimmten Barometerstand nach jener Hypothese eine größere Höhe, während seine Rechnung eine kleinere ergibt. Für den Monte Gregorio gibt er für diesen Unterschied $-16$ Fuß an. Richtig sind $+2$ Fuß. Bei kleineren Höhen ist der Unterschied sehr nahe proportional zum Quadrat der Höhe. Das ließ sich auch durch eine naheliegende Überlegung leicht \textit{a priori} voraussehen, die wir hier jedoch der Kürze halber nicht weiter ausführen. Die Unterschiede des Herrn \textsc{Benzenberg} sind dagegen bei kleinen Höhen nahezu proportional zu diesen selbst. Schon das genügte, um ihre Unrichtigkeit zu erkennen. Im Übrigen ist die obige numerische Rechnung nur für trockene Luft ausgeführt. Die Berücksichtigung des Wasserdampfs würde die relativen Unterschiede nur sehr wenig ändern.

Wir fügen dieser Darstellung noch einige Bemerkungen hinzu.

I.\quad Unser Ergebnis besagt, dass der Unterschied zwischen der Barometerhöhe nach \textsc{Dalton}s Hypothese und der auf gewöhnliche Weise berechneten \textit{positiv} ist und bei mäßigen Höhen annähernd proportional zu ihren Quadraten wird. Es gilt allgemein und hängt von den angenommenen Werten der spezifischen Gewichte der einzelnen Gase, aus denen die gemischte Luft besteht, nicht ab. Es wäre also vergeblich, von anderen Werten dieser spezifischen Gewichte ein günstigeres Ergebnis zu erwarten.

II.\quad Schon 1807 hat \textsc{Tralles} richtig dargestellt, wie man die Barometerhöhen nach \textsc{Dalton}s Hypothese berechnet. Schon ein flüchtiger Blick darauf genügt, um zu erkennen, dass sein Ergebnis im Wesentlichen ganz mit unserem übereinstimmt. Man muss sich daher wundern, dass der Verfasser des einen oben erwähnten Artikels des physikalischen Wörterbuchs behauptet, \textsc{Tralles} finde fast genau dieselbe Differenz wie Herr \textsc{Benzenberg}. Tatsächlich unterscheidet sie sich sowohl im Vorzeichen als auch in der Gesetzmäßigkeit. Bei mäßigen
\SourceLeaf{596}Höhen ist sie proportional zum Quadrat der Höhe. Zur Höhe selbst ist sie nicht proportional. Vielleicht ist ein Druckfehler, der allerdings leicht als solcher zu erkennen ist, an diesem Versehen schuld. Denn im Aufsatz von \textsc{Tralles} (\textsc{Gilbert}s Annalen, Bd.\ 27, S.\ 445) steht einmal $\frac{B}{b+\beta}$ anstelle von $\frac{b+\beta}{B}$. Bei der Anwendung auf ein bestimmtes Beispiel (ebenda, unterste Zeile) steht der Bruch aber wieder richtig, ebenso ist auf der folgenden Seite der Unterschied mit dem richtigen Vorzeichen angesetzt. Noch mehr muss man sich aber wundern, dass Herr \textsc{Benzenberg} es unterließ, das Ergebnis seines als gründlicher Mathematiker bekannten Vorgängers mit seinem eigenen zu vergleichen, dem es ganz entgegengesetzt ist. Man würde in der Tat vermuten, dass Herr \textsc{Benzenberg} diesen Aufsatz gar nicht gekannt hat, wenn er ihn nicht ausdrücklich erwähnte, wenn auch nur mit der abfertigenden Bemerkung auf S.\ 15. „Herr \textsc{Tralles} hat Buchstabenrechnung verwendet. Das ist unnötig. Wenn man sich vier Barometer vorstellt, kann man es mit dem Dreisatz ausführen, und man braucht überhaupt keine Gelehrsamkeit“. Diesen Grundsatz, zu dem sich Herr \textsc{Benzenberg} bei vielen Gelegenheiten laut bekannt hat — ob immer bei den richtigen, wollen wir hier nicht untersuchen —, mag man im vorliegenden Fall immerhin gelten lassen. Unsere Darstellung mag dann selbst als Bestätigung dienen, wenn es uns gelungen ist, ihr die nötige Klarheit zu geben.

III.\quad Der Versuch, den Unterschied zwischen der barometrischen und der trigonometrischen Messung des Monte Gregorio durch \textsc{Dalton}s Hypothese zu beseitigen, ist vergeblich (nach ihr wird der Unterschied sogar noch um zwei Fuß größer). Damit steht als gesicherte Tatsache fest, dass eine der beiden Messungen oder beide nicht die Genauigkeit haben, die Herr \textsc{Benzenberg} ihnen zuschreiben zu können glaubte. Unserer Meinung nach mögen alle drei hier infrage kommenden Fehlerquellen daran beteiligt sein. Erstens die Schwankungen der gemessenen Barometerhöhen selbst. Zweitens die in der Rechnung verwendete Konstante, die Herr \textsc{Benzenberg} auf \textsc{Biot}s Wägung der atmosphärischen Luft gestützt hatte und die sich wohl viel zuverlässiger bestimmen lässt, wenn man zahlreiche sowohl barometrisch als auch trigonometrisch gemessene Berghöhen zweckmäßig heranzieht. Drittens mag auch die trigonometrische Messung des Monte Gregorio selbst zu dem Unterschied beigetragen haben, da wir dem Urteil des Herrn \textsc{Benzenberg} über ihre unübertreffliche Genauigkeit nicht ganz zustimmen können. Die Bestimmung der Ent\SourceLeaf{597}fernung zwischen der Bergspitze und dem Standpunkt am Fuß des Berges beruhte nur auf einer kurzen Basis. Ihre Länge (670 Meter) scheint zwar mit großer Sorgfalt gemessen worden zu sein (man verbrachte vier Tage damit). Der ihr gegenüberliegende Winkel (nur $6^0\;14'$) wurde aber nur aus den beiden anderen erschlossen und nicht selbst gemessen. Ein solches Verfahren erfordert selbst beim Einsatz hervorragender Instrumente stets große Vorsicht. Doch \textsc{Daubuisson}s Winkelmessinstrument, ein \textsc{Lenoir}scher Repetitionskreis von acht Zoll Durchmesser, scheint sehr mittelmäßig gewesen zu sein. Denn von den zehn Wiederholungen, aus denen die Winkelmessungen an jedem der beiden Standpunkte bestehen, liefern die einzelnen Paare einmal Unterschiede für den einfachen Winkel, die über eine Minute hinausgehen. Auch sagt \textsc{Daubuisson} nichts über die Beschaffenheit der Signale, die er als Zielpunkte benutzte. Daher lässt sich nicht beurteilen, mit welcher Schärfe sie sich einschneiden ließen und ob dabei womöglich eine ungünstige Beleuchtungsphase eintreten konnte. Ein Fehler von einer halben Minute im erschlossenen dritten Winkel würde die gemessene Höhe schon um sieben Fuß ändern. — Im Übrigen können wir es dabei belassen, denn für die gegenwärtige Frage ist es ganz gleichgültig, wie man den Unterschied erklären will. Es genügt, dass \textsc{Dalton}s Hypothese gar nichts zu seiner Erklärung beitragen kann.

